\subsection{Example: Representing Sets}
\label{Section 2.3.3}

In the previous examples we built representations for two kinds of compound
data objects: rational numbers and algebraic expressions.  In one of these
examples we had the choice of simplifying (reducing) the expressions at either
construction time or selection time, but other than that the choice of a
representation for these structures in terms of lists was straightforward. When
we turn to the representation of sets, the choice of a representation is not so
obvious.  Indeed, there are a number of possible representations, and they
differ significantly from one another in several ways.

Informally, a set is simply a collection of distinct objects.  To give a more
precise definition we can employ the method of data abstraction.  That is, we
define ``set'' by specifying the operations that are to be used on sets.  These
are \code{union\_set}, \code{intersection\_set}, \code{is\_element\_of\_set},
and \code{adjoin\_set}.  \code{is\_element\_of\_set} is a predicate that
determines whether a given element is a member of a set.  \code{adjoin\_set}
takes an object and a set as arguments and returns a set that contains the
elements of the original set and also the adjoined element.  \code{union\_set}
computes the union of two sets, which is the set containing each element that
appears in either argument.  \code{intersection\_set} computes the intersection
of two sets, which is the set containing only elements that appear in both
arguments. From the viewpoint of data abstraction, we are free to design any
representation that implements these operations in a way consistent with the
interpretations given above.\footnote{If we want to be more formal, we can
specify ``consistent with the interpretations given above'' to mean that the
operations satisfy a collection of rules such as these:

\noindent
\( \bullet \) For any set \code{S} and any object \code{x},
\code{is\_element\_of\_set(x, adjoin\_set(x, S))}
is true (informally: ``Adjoining an object to a
set produces a set that contains the object'').

\noindent
\( \bullet \) For any sets \code{S} and \code{T} and any object \code{x},
\code{is\_element\_of\_set(x, union\_set(S, T))}
is equal to
\code{is\_element\_of\_set(x, S) or is\_element\_of\_set(x, T)}
(informally: ``The elements of \code{union\_set(S, T)} are the elements that
are in \code{S} or in \code{T}'').

\noindent
\( \bullet \) For any object \code{x},
\code{is\_element\_of\_set(x, ())}
is false (informally: ``No object is an element of the empty set'').
}

\begin{revealTheSurface}
Python has a set.  It is built into the language, and it does every operation
named above---membership, adjoining, union, intersection---faster than anything
this section will build, including the binary tree it ends with.  Saying so at
the outset seems the honest course, since the reader will find out within the
hour in any case.

Two of them, in fact.  Written with braces, \code{\{1, 3, 5\}} is a \code{set},
and a \code{set} can be changed after it is made: elements may be added to it
and taken out of it.  \code{frozenset(\{1, 3, 5\})} is the same thing with that
ability removed, and it is the one we shall use, for the reason given at the
opening of this chapter---the original's data does not change once made, and
where we write in the Pythonic manner we keep to that.  The mutable one is not
a mistake and we are not avoiding it out of superstition; it is the subject of
\link{Chapter 3}, and we shall take it up when the book does.

That does not make the section idle, and it is worth being clear about what it
is now for.  It is the book's most sustained demonstration that a
representation is a \emph{choice}: the same four operations are implemented
three times, the interface never changes, and the cost changes by orders of
magnitude.  The programs that use a set are written once and do not care which
of the three they are handed.  That is the lesson, and having a fast set in the
language does nothing to it---one still has to know why a tree beats a list, if
only to understand what the built-in set is doing instead and when it stops
being the right answer.

A word, too, about what the exercise will cost us.  As everywhere in this
chapter we build on tuples, and our data does not change once it is made:
\code{adjoin\_set} does not put an element into a set, it constructs a new set
with the element in it, and the procedures that walk a set do so by taking its
rest---\code{s[1:]}---rather than by following a pointer.  Both operations
copy, in the way \link{Section 2.2.3} set out, where the original's cost
nothing.

The consequence is that the order-of-growth figures quoted in this section are
the \emph{original's}, and ours are worse.  Where the original's scan of an
unordered list is \( \Theta(n) \), ours is \( \Theta(n^2) \), because each of
the \( n \) steps copies what remains of the tuple.  The tree escapes this, and
it is worth seeing why: its selectors index into a tuple---\code{s[1]},
\code{s[2]}---rather than slicing it, and indexing costs nothing.  Our tree
really is \( \Theta(\log n) \); sixty-four times the elements cost about half
again the time.  So the gap between the list and the tree is \emph{wider} for
us than for the original, and the section's argument comes out stronger rather
than weaker---but the figures in the text should be read as statements about
Scheme, and re-derived where they matter.

We write it this way for consistency with the rest of the chapter, and because
the shape of the three representations is what is being taught.  Python's own
set is not built this way at all, as the end of the section will show.

We shall build all three, as the original does, and come back to Python's own
set then, when there is something to compare it against.
\end{revealTheSurface}

\subsubsection*{Sets as unordered lists}

One way to represent a set is as a list of its elements in which no element
appears more than once.  The empty set is represented by the empty list.  In
this representation, \code{is\_element\_of\_set} is the procedure \code{memq}
of \link{Section 2.3.1} under another name.\footnote{The original has to say
more here.  It has two notions of sameness---\code{eq?}, which asks whether two
symbols are the same object, and \code{equal?}, which compares lists element by
element---and must be careful to use the second, so that the elements of a set
need not be symbols.  We have only \code{==}, which compares numbers by value
and sequences element by element and recursively, so the choice does not arise
and \code{memq} was already doing the more general thing.}

\begin{codeBlock}
>>> def is_element_of_set(x, s):
...     if s == ():
...         return False
...     elif x == s[0]:
...         return True
...     else:
...         return is_element_of_set(x, s[1:])
\end{codeBlock}

\noindent
Using this, we can write \code{adjoin\_set}.  If the object to be adjoined is
already in the set, we just return the set.  Otherwise, we use \code{cons} to
add the object to the list that represents the set:

\begin{codeBlock}
>>> def adjoin_set(x, s):
...     if is_element_of_set(x, s):
...         return s
...     else:
...         return (x,) + s
\end{codeBlock}

\noindent
For \code{intersection\_set} we can use a recursive strategy.  If we know how
to form the intersection of \code{set2} and the \code{cdr} of \code{set1}, we
only need to decide whether to include the \code{car} of \code{set1} in this.
But this depends on whether \code{(car set1)} is also in \code{set2}.  Here is
the resulting procedure:

\begin{codeBlock}
>>> def intersection_set(s1, s2):
...     if s1 == () or s2 == ():
...         return ()
...     elif is_element_of_set(s1[0], s2):
...         return (s1[0],) + intersection_set(s1[1:], s2)
...     else:
...         return intersection_set(s1[1:], s2)
\end{codeBlock}

\noindent
In designing a representation, one of the issues we should be concerned with is
efficiency.  Consider the number of steps required by our set operations. Since
they all use \code{is\_element\_of\_set}, the speed of this operation has a
major impact on the efficiency of the set implementation as a whole.  Now, in
order to check whether an object is a member of a set,
\code{is\_element\_of\_set} may have to scan the entire set. (In the worst
case, the object turns out not to be in the set.)  Hence, if the set has \( n
\) elements, \code{is\_element\_of\_set}  might take up to \( n \) steps.
Thus, the number of steps required grows as \( \Theta(n) \).  The number of
steps required by \code{adjoin\_set}, which uses this operation, also grows as
\( \Theta(n) \). For \code{intersection\_set}, which does an
\code{is\_element\_of\_set} check for each element of \code{set1}, the number
of steps required grows as the product of the sizes of the sets involved, or \(
\Theta(n^2) \) for two sets of size \( n \).  The same will be true of
\code{union\_set}.

\begin{exerciseGroup}
\phantomsection\label{Exercise 2.59}
\exerciseHeading{Exercise 2.59:} Implement the \code{union\_set}
operation for the unordered-list representation of sets.

\phantomsection\label{Exercise 2.60}
\exerciseHeading{Exercise 2.60:} We specified that a set would be
represented as a list with no duplicates.  Now suppose we allow duplicates. For
instance, the set \( \{1, 2, 3\} \) could be represented as the sequence
\code{(2, 3, 2, 1, 3, 2, 2)}.  Design procedures \code{is\_element\_of\_set},
\code{adjoin\_set}, \code{union\_set}, and \code{intersection\_set} that
operate on this representation.  How does the efficiency of each compare with
the corresponding procedure for the non-duplicate representation?  Are there
applications for which you would use this representation in preference to the
non-duplicate one?
\end{exerciseGroup}

\subsubsection*{Sets as ordered lists}

One way to speed up our set operations is to change the representation so that
the set elements are listed in increasing order.  To do this, we need some way
to compare two objects so that we can say which is bigger.  For example, we
could compare symbols lexicographically, or we could agree on some method for
assigning a unique number to an object and then compare the elements by
comparing the corresponding numbers.  To keep our discussion simple, we will
consider only the case where the set elements are numbers, so that we can
compare elements using \code{>} and \code{<}.  We will represent a set of
numbers by listing its elements in increasing order.  Whereas our first
representation above allowed us to represent the set \( \{1, 3, 6, 10\} \) by
listing the elements in any order, our new representation allows only the list
\code{(1, 3, 6, 10)}.

One advantage of ordering shows up in \code{is\_element\_of\_set}: In checking
for the presence of an item, we no longer have to scan the entire set.  If we
reach a set element that is larger than the item we are looking for, then we
know that the item is not in the set:

\begin{codeBlock}
>>> def is_element_of_set(x, s):
...     if s == ():
...         return False
...     elif x == s[0]:
...         return True
...     elif x < s[0]:
...         return False
...     else:
...         return is_element_of_set(x, s[1:])
\end{codeBlock}

\noindent
How many steps does this save?  In the worst case, the item we are looking for
may be the largest one in the set, so the number of steps is the same as for
the unordered representation.  On the other hand, if we search for items of
many different sizes we can expect that sometimes we will be able to stop
searching at a point near the beginning of the list and that other times we
will still need to examine most of the list.  On the average we should expect
to have to examine about half of the items in the set.  Thus, the average
number of steps required will be about \( n / 2 \).  This is still
\( \Theta(n) \) growth, but it does save us, on the average, a factor of 2
in number of steps over the previous implementation.

We obtain a more impressive speedup with \code{intersection\_set}.  In the
unordered representation this operation required \( \Theta(n^2) \) steps,
because we performed a complete scan of \code{set2} for each element of
\code{set1}.  But with the ordered representation, we can use a more clever
method.  Begin by comparing the initial elements, \code{x1} and \code{x2}, of
the two sets.  If \code{x1} equals \code{x2}, then that gives an element of the
intersection, and the rest of the intersection is the intersection of the
\code{cdr}-s of the two sets.  Suppose, however, that \code{x1} is less than
\code{x2}.  Since \code{x2} is the smallest element in \code{set2}, we can
immediately conclude that \code{x1} cannot appear anywhere in \code{set2} and
hence is not in the intersection.  Hence, the intersection is equal to the
intersection of \code{set2} with the \code{cdr} of \code{set1}.  Similarly, if
\code{x2} is less than \code{x1}, then the intersection is given by the
intersection of \code{set1} with the \code{cdr} of \code{set2}.  Here is the
procedure:

\begin{codeBlock}
>>> def intersection_set(s1, s2):
...     if s1 == () or s2 == ():
...         return ()
...     x1, x2 = s1[0], s2[0]
...     if x1 == x2:
...         return (x1,) + intersection_set(s1[1:], s2[1:])
...     elif x1 < x2:
...         return intersection_set(s1[1:], s2)
...     else:
...         return intersection_set(s1, s2[1:])
\end{codeBlock}

\noindent
To estimate the number of steps required by this process, observe that at each
step we reduce the intersection problem to computing intersections of smaller
sets---removing the first element from \code{set1} or \code{set2} or both.
Thus, the number of steps required is at most the sum of the sizes of
\code{set1} and \code{set2}, rather than the product of the sizes as with the
unordered representation.  This is \( \Theta(n) \) growth rather than
\( \Theta(n^2) \)---a considerable speedup, even for sets of moderate size.

\begin{exerciseGroup}
\phantomsection\label{Exercise 2.61}
\exerciseHeading{Exercise 2.61:} Give an implementation of
\code{adjoin\_set} using the ordered representation.  By analogy with
\code{is\_element\_of\_set} show how to take advantage of the ordering to
produce a
procedure that requires on the average about half as many steps as with the
unordered representation.

\phantomsection\label{Exercise 2.62}
\exerciseHeading{Exercise 2.62:} Give a \( \Theta(n) \)
implementation of \code{union\_set} for sets represented as ordered lists.
\end{exerciseGroup}

\subsubsection*{Sets as binary trees}

We can do better than the ordered-list representation by arranging the set
elements in the form of a tree.  Each node of the tree holds one element of the
set, called the ``entry'' at that node, and a link to each of two other
(possibly empty) nodes.  The ``left'' link points to elements smaller than the
one at the node, and the ``right'' link to elements greater than the one at the
node.  \link{Figure 2.15} shows some trees that represent the set \( \{1, 3, 5,
7, 9, 11\} \).  The same set may be represented by a tree in a number of
different ways.  The only thing we require for a valid representation is that
all elements in the left subtree be smaller than the node entry and that all
elements in the right subtree be larger.

% figure 2.16
\begin{imageFigure}{Various binary trees that represent the set
  \( \{1, 3, 5, 7, 9, 11\} \).}
  \resizebox{\linewidth}{!}{% Three binary trees that all represent the set {1, 3, 5, 7, 9, 11}, for
% Section 2.3.3. Redrawn from upstream's figure 2.16.
%
% Each is given as its entry and the positions of its nodes, so the three can
% share one drawing routine: a node at (x, y) with an edge up to its parent.
\begin{tikzpicture}[
  x = 1mm, y = -1mm,
  every node/.style = {font=\ttfamily\footnotesize, inner sep=1.2pt},
]

% left tree: 7 at the root
\foreach \v/\x/\y in {7/20/0, 3/10/14, 9/30/14, 1/4/28, 5/16/28, 11/26/28}
  \node (a\v) at (\x, \y) {\v};
\foreach \p/\c in {a7/a3, a7/a9, a3/a1, a3/a5, a9/a11}
  \draw[thin] (\p.south) -- (\c.north);

% middle tree: 3 at the root, leaning right
\begin{scope}[xshift=50mm]
\foreach \v/\x/\y in {3/16/0, 1/6/14, 7/26/14, 5/20/28, 9/34/28, 11/40/42}
  \node (b\v) at (\x, \y) {\v};
\foreach \p/\c in {b3/b1, b3/b7, b7/b5, b7/b9, b9/b11}
  \draw[thin] (\p.south) -- (\c.north);
\end{scope}

% right tree: 5 at the root, balanced
\begin{scope}[xshift=104mm]
\foreach \v/\x/\y in {5/20/0, 3/10/14, 9/30/14, 1/4/28, 7/24/28, 11/38/28}
  \node (c\v) at (\x, \y) {\v};
\foreach \p/\c in {c5/c3, c5/c9, c3/c1, c9/c7, c9/c11}
  \draw[thin] (\p.south) -- (\c.north);
\end{scope}

\end{tikzpicture}
}
\end{imageFigure}

The advantage of the tree representation is this: Suppose we want to check
whether a number \( x \) is contained in a set.  We begin by comparing \( x \)
with the entry in the top node.  If \( x \) is less than this, we know that we
need only search the left subtree; if \( x \) is greater, we need only search
the right subtree.  Now, if the tree is ``balanced,'' each of these subtrees
will be about half the size of the original.  Thus, in one step we have reduced
the problem of searching a tree of size \( n \) to searching a tree of size \(
n / 2 \). Since the size of the tree is halved at each step, we should expect
that the number of steps needed to search a tree of size \( n \) grows as \(
\Theta(\log n) \).\footnote{Halving the size of the problem at each step is the
distinguishing characteristic of logarithmic growth, as we saw with the
fast-exponentiation algorithm of \link{Section 1.2.4} and the half-interval
search method of \link{Section 1.3.3}.} For large sets, this will be a
significant speedup over the previous representations.

We can represent trees by using lists.  Each node will be a list of three
items: the entry at the node, the left subtree, and the right subtree.  A left
or a right subtree of the empty list will indicate that there is no subtree
connected there.  We can describe this representation by the following
procedures:\footnote{We are representing sets in terms of trees, and trees in
terms of lists---in effect, a data abstraction built upon a data abstraction.
We can regard the procedures \code{entry}, \code{left\_branch},
\code{right\_branch}, and \code{make\_tree} as a way of isolating the
abstraction of a ``binary tree'' from the particular way we might wish to
represent such a tree in terms of list structure.}

\begin{codeBlock}
>>> def entry(tree):
...     return tree[0]

>>> def left_branch(tree):
...     return tree[1]

>>> def right_branch(tree):
...     return tree[2]

>>> def make_tree(entry, left, right):
...     return (entry, left, right)
\end{codeBlock}

\noindent
Now we can write the \code{is\_element\_of\_set} procedure using the strategy
described above:

\begin{codeBlock}
>>> def is_element_of_set(x, s):
...     if s == ():
...         return False
...     elif x == entry(s):
...         return True
...     elif x < entry(s):
...         return is_element_of_set(x, left_branch(s))
...     else:
...         return is_element_of_set(x, right_branch(s))
\end{codeBlock}

\noindent
Adjoining an item to a set is implemented similarly and also requires
\( \Theta(\log n) \) steps.  To adjoin an item \code{x}, we compare
\code{x} with the node entry to determine whether \code{x} should be added to
the right or to the left branch, and having adjoined \code{x} to the
appropriate branch we piece this newly constructed branch together with the
original entry and the other branch.  If \code{x} is equal to the entry, we
just return the node.  If we are asked to adjoin \code{x} to an empty tree, we
generate a tree that has \code{x} as the entry and empty right and left
branches.  Here is the procedure:

\begin{codeBlock}
>>> def adjoin_set(x, s):
...     if s == ():
...         return make_tree(x, (), ())
...     elif x == entry(s):
...         return s
...     elif x < entry(s):
...         return make_tree(entry(s),
...                          adjoin_set(x, left_branch(s)),
...                          right_branch(s))
...     else:
...         return make_tree(entry(s),
...                          left_branch(s),
...                          adjoin_set(x, right_branch(s)))
\end{codeBlock}

% figure 2.17
\begin{imageFigure}{Unbalanced tree produced by adjoining 1 through 7 in
  sequence.}
  % The tree produced by adjoining 1 through 7 in sequence, for Section 2.3.3.
% Redrawn from upstream's figure 2.17. Every left branch is empty, so the tree
% is a list in disguise and searching it is no faster than scanning one.
\begin{tikzpicture}[
  x = 1mm, y = -1mm,
  every node/.style = {font=\ttfamily\footnotesize, inner sep=1.2pt},
]

\foreach \v/\i in {1/0, 2/1, 3/2, 4/3, 5/4, 6/5, 7/6}
  \node (n\v) at (\i*7, \i*11) {\v};
\foreach \p/\c in {n1/n2, n2/n3, n3/n4, n4/n5, n5/n6, n6/n7}
  \draw[thin] (\p.south) -- (\c.north);

\end{tikzpicture}

\end{imageFigure}

\noindent
The above claim that searching the tree can be performed in a logarithmic
number of steps rests on the assumption that the tree is ``balanced,'' i.e.,
that the left and the right subtree of every tree have approximately the same
number of elements, so that each subtree contains about half the elements of
its parent.  But how can we be certain that the trees we construct will be
balanced?  Even if we start with a balanced tree, adding elements with
\code{adjoin\_set} may produce an unbalanced result.  Since the position of a
newly adjoined element depends on how the element compares with the items
already in the set, we can expect that if we add elements ``randomly'' the tree
will tend to be balanced on the average.  But this is not a guarantee.  For
example, if we start with an empty set and adjoin the numbers 1 through 7 in
sequence we end up with the highly unbalanced tree shown in \link{Figure 2.16}.
In this tree all the left subtrees are empty, so it has no advantage over a
simple ordered list.  One way to solve this problem is to define an operation
that transforms an arbitrary tree into a balanced tree with the same elements.
Then we can perform this transformation after every few \code{adjoin\_set}
operations to keep our set in balance.  There are also other ways to solve this
problem, most of which involve designing new data structures for which
searching and insertion both can be done in \( \Theta(\log n) \)
steps.\footnote{Examples of such structures include \newterm{B-trees} and
\newterm{red-black trees}.  There is a large literature on data structures
devoted to this problem.  See \link{Cormen et al. 1990}.}

\begin{exerciseGroup}
\phantomsection\label{Exercise 2.63}
\exerciseHeading{Exercise 2.63:} Each of the following two
procedures converts a binary tree to a list.

\begin{codeBlock}
>>> def tree_to_list_1(tree):
...     if tree == ():
...         return ()
...     else:
...         return (tree_to_list_1(left_branch(tree))
...                 + (entry(tree),)
...                 + tree_to_list_1(right_branch(tree)))

>>> def tree_to_list_2(tree):
...     def copy_to_list(tree, result):
...         if tree == ():
...             return result
...         else:
...             return copy_to_list(
...                 left_branch(tree),
...                 (entry(tree),) + copy_to_list(
...                     right_branch(tree), result))
...     return copy_to_list(tree, ())
\end{codeBlock}

\begin{enumerate}[a.]

\item
Do the two procedures produce the same result for every tree?  If not, how do
the results differ?  What lists do the two procedures produce for the trees in
\link{Figure 2.15}?

\item
Do the two procedures have the same order of growth in the number of steps
required to convert a balanced tree with \( n \) elements to a list?  If not,
which one grows more slowly?

\end{enumerate}

\phantomsection\label{Exercise 2.64}
\exerciseHeading{Exercise 2.64:} The following procedure
\code{list\_to\_tree} converts an ordered list to a balanced binary tree.  The
helper procedure \code{partial\_tree} takes as arguments an integer \( n \) and
list of at least \( n \) elements and constructs a balanced tree containing the
first \( n \) elements of the list.  The result returned by
\code{partial\_tree} is a pair whose first element is the constructed tree and
whose second is the sequence of elements not included in it.\footnote{The
original returns this pair and takes it apart with \code{car} and \code{cdr};
we unpack it into two names on the left of the assignment, which Python allows
for any sequence of the right length.  It is the same pair, read more
directly.}

\begin{codeBlock}
>>> def list_to_tree(elements):
...     return partial_tree(elements, len(elements))[0]

>>> def partial_tree(elts, n):
...     if n == 0:
...         return ((), elts)
...     left_size = (n - 1) // 2
...     left_tree, non_left_elts = partial_tree(elts, left_size)
...     this_entry = non_left_elts[0]
...     right_size = n - left_size - 1
...     right_tree, remaining = partial_tree(
...         non_left_elts[1:], right_size)
...     return (make_tree(this_entry, left_tree, right_tree),
...             remaining)
\end{codeBlock}

\begin{enumerate}[a.]

\item
Write a short paragraph explaining as clearly as you can how
\code{partial\_tree} works.  Draw the tree produced by \code{list\_to\_tree}
for the sequence \code{(1, 3, 5, 7, 9, 11)}.

\item
What is the order of growth in the number of steps required by
\code{list\_to\_tree} to convert a list of \( n \) elements?

\end{enumerate}

\phantomsection\label{Exercise 2.65}
\exerciseHeading{Exercise 2.65:} Use the results of \link{Exercise 2.63}
and \link{Exercise 2.64} to give \( \Theta(n) \) implementations of
\code{union\_set} and \code{intersection\_set} for sets implemented as
(balanced) binary trees.\footnote{\link{Exercise 2.63} through \link{Exercise
2.65} are due to Paul Hilfinger.}
\end{exerciseGroup}

\subsubsection*{Sets and information retrieval}

We have examined options for using lists to represent sets and have seen how
the choice of representation for a data object can have a large impact on the
performance of the programs that use the data.  Another reason for
concentrating on sets is that the techniques discussed here appear again and
again in applications involving information retrieval.

Consider a data base holding a large number of individual records, such as the
personnel files for a company, or the transactions in an accounting system. A
typical data-management system spends a large amount of time accessing or
modifying the data in the records and therefore requires an efficient method
for accessing records.  This is done by identifying a part of each record to
serve as an identifying \newterm{key}.  A key can be anything that uniquely
identifies the record.  For a personnel file, it might be an employee's
\acronym{ID} number.  For an accounting system, it might be a transaction
number.  Whatever the key is, when we define the record as a data structure we
should include a \code{key} selector procedure that retrieves the key
associated with a given record.

Now we represent the data base as a set of records. To locate the record with a
given key we use a procedure \code{lookup}, which takes as arguments a key and
a data base and which returns the record that has that key, or false if there
is no such record.  \code{lookup} is implemented in almost the same way as
\code{is\_element\_of\_set}.  For example, if the set of records is implemented
as an unordered list, we could use

\begin{codeBlock}
>>> def lookup(given_key, records):
...     if records == ():
...         return False
...     elif given_key == key(records[0]):
...         return records[0]
...     else:
...         return lookup(given_key, records[1:])
\end{codeBlock}

\noindent
Of course, there are better ways to represent large sets than as unordered
lists.  Information-retrieval systems in which records have to be ``randomly
accessed'' are typically implemented by a tree-based method, such as the
binary-tree representation discussed previously.  In designing such a system
the methodology of data abstraction can be a great help.  The designer can
create an initial implementation using a simple, straightforward representation
such as unordered lists.  This will be unsuitable for the eventual system, but
it can be useful in providing a ``quick and dirty'' data base with which to
test the rest of the system.  Later on, the data representation can be modified
to be more sophisticated.  If the data base is accessed in terms of abstract
selectors and constructors, this change in representation will not require any
changes to the rest of the system.

\begin{exerciseGroup}
\phantomsection\label{Exercise 2.66}
\exerciseHeading{Exercise 2.66:} Implement the \code{lookup}
procedure for the case where the set of records is structured as a binary tree,
ordered by the numerical values of the keys.
\end{exerciseGroup}


\subsubsection*{The set Python already has}

\begin{revealTheSurface}
We can now say what we deferred at the start of the section.  The four
operations are an operator apiece:

\begin{codeBlock}
>>> s = frozenset({1, 3, 5, 7, 9, 11})
>>> 5 in s
True
>>> 4 in s
False
>>> s | frozenset({2})
frozenset({1, 2, 3, 5, 7, 9, 11})
>>> s & frozenset({3, 4, 5})
frozenset({3, 5})
\end{codeBlock}

\noindent
\code{in} is \code{is\_element\_of\_set}, \code{|} is \code{union\_set},
\code{\&} is \code{intersection\_set}, and adjoining is
\code{s | frozenset(\{x\})}.  So far this is only notation.  The interesting
part is the cost.  Membership in a frozenset of two hundred thousand elements,
a thousand times over:

\begin{codeBlock}
0.44 milliseconds
\end{codeBlock}

\noindent
That is not \( \Theta(\log n) \).  It is not proportional to anything: doubling
the size of the set would not change it.  Membership in a Python set takes
constant time on average, which beats the best of the three representations
this section built, and by a margin that grows without limit as the set grows.

It does this by \newterm{hashing}.  Instead of comparing the sought element
with the elements of the set, Python computes a number from the element
itself---its \newterm{hash}---and uses that number to decide where the element
would have to be stored if it were present.  Then it looks only there.  The
search does not get longer as the set gets larger, because it was never a
search: nothing is compared except the handful of elements that happen to hash
to the same place.

What this costs is the ordering.  Our tree keeps its elements sorted, and can
answer ``what is the smallest element greater than \( x \)'' in
\( \Theta(\log n) \); a hash set cannot answer it at all without looking at
everything, because the hash of an element tells you nothing about the hash of
its neighbours.  The original's note that information-retrieval systems are
``typically implemented by a tree-based method'' has not stopped being true.
Trees are what a database index is made of, and the reason is exactly this: a
hash table finds one record faster, and a tree finds a \emph{range} of them.

There is one requirement.  To be put in a set, an object must be hashable, and
must promise that two objects which compare equal hash alike---otherwise the
same element could be sought in one place and stored in another.  Python
supplies a \code{\_\_hash\_\_} for its own types and for any class that does
not define \code{\_\_eq\_\_}.  Define \code{\_\_eq\_\_}, as our \code{Symbol}
of \link{Section 2.3.2} does, and Python removes the inherited
\code{\_\_hash\_\_}---on the reasoning that a class which has redefined
equality has almost certainly invalidated the old hash:

\begin{codeBlock}
>>> x = Symbol("x")
>>> frozenset({x})
TypeError: cannot use 'Symbol' as a set element
           (unhashable type: 'Symbol')
\end{codeBlock}

\noindent
The remedy is to supply one, hashing on exactly the thing equality compares:

\begin{codeBlock}
>>> class Symbol(Expression):
...     def __init__(self, name):
...         self.name = name
...     def __repr__(self):
...         return self.name
...     def __eq__(self, other):
...         return isinstance(other, Symbol) and self.name == other.name
...     def __hash__(self):
...         return hash(self.name)

>>> x, y = Symbol("x"), Symbol("y")
>>> frozenset({x, y, Symbol("x")})
frozenset({x, y})
\end{codeBlock}

\noindent
Note the last line: three symbols went in and two came out, because the second
\code{x} compared equal to the first and hashed alike, so the set kept one of
them.  That is the definition of a set doing its work, and it works only
because the two methods agree.  A \code{\_\_hash\_\_} that disagreed with
\code{\_\_eq\_\_} would not raise an error; it would quietly produce a set
containing the same element twice.
\end{revealTheSurface}
